\section{Softwarearchitektur}
Die Softwarearchitektur der Wortuhr îst als modulare Struktur aufgesetzt, in der einzelne Funktionsblöcke klar voneinander getrennt sind. 
Die Bestandteile sind die Initialisierung des Systems, die zyklische Ablaufsteuerung, die Zeitverarbeitung, die Anzeige-Logik sowie die Zusatzfunktionen durch Sensorik und Externe \gls{api}-Schnittstellen.
Die Kommunikation zwischen den Komponenten erfolgts über globale Zustandsvariablen, welche in festen Zeitintervallen aktualisiert und anschließend zur Anzeige verarbeitet werden.

\subsection{Initialisierung des Systems}

Die Initialisierung erfolgt einmalig beim Systemstart. 
Dabei werden alle Hardwarekomponenten konfiguriert, die Netzwerkverbindung aufgebaut sowie die aktuelle Zeit über einen \gls{ntp}-Server synchronisiert. 
Zusätzlich werden erste Sensorwerte und Wetterdaten geladen, um direkt nach dem Start eine gültige Anzeige zu gewährleisten. 
Der Ablauf der Initialisierung ist in Abbildung~\ref{fig:setup} dargestellt.

\begin{figure}[H]
\centering
\begin{tikzpicture}[
    node distance=1.1cm,
    startstop/.style={rectangle, rounded corners, minimum width=4cm, minimum height=0.8cm, draw=black},
    process/.style={rectangle, minimum width=4.8cm, minimum height=0.8cm, draw=black},
    arrow/.style={->, thick}
]

\node (start) [startstop] {Start};
\node (s1) [process, below=of start] {Serielle Schnittstelle starten};
\node (s2) [process, below=of s1] {LED-Streifen initialisieren};
\node (s3) [process, below=of s2] {Sensor initialisieren};
\node (s4) [process, below=of s3] {WLAN verbinden};
\node (s5) [process, below=of s4] {NTP konfigurieren};
\node (s6) [process, below=of s5] {Zeit synchronisieren};
\node (s7) [process, below=of s6] {Wetterdaten abrufen};
\node (s8) [process, below=of s7] {Sensorwerte einlesen};
\node (s9) [process, below=of s8] {Erste Anzeige erzeugen};
\node (s10) [process, below=of s9] {HTTP-Server starten};
\node (end) [startstop, below=of s10] {System bereit};

\draw [arrow] (start) -- (s1);
\draw [arrow] (s1) -- (s2);
\draw [arrow] (s2) -- (s3);
\draw [arrow] (s3) -- (s4);
\draw [arrow] (s4) -- (s5);
\draw [arrow] (s5) -- (s6);
\draw [arrow] (s6) -- (s7);
\draw [arrow] (s7) -- (s8);
\draw [arrow] (s8) -- (s9);
\draw [arrow] (s9) -- (s10);
\draw [arrow] (s10) -- (end);

\end{tikzpicture}
\caption{Ablauf der Systeminitialisierung}
\label{fig:setup}
\end{figure}

\subsection{Zyklische Ablaufsteuerung}

Nach der Initialisierung wird das System in einer Endlosschleife betrieben. 
Innerhalb dieser Schleife werden zyklisch Webanfragen verarbeitet, Zeitinformationen aktualisiert sowie Anzeige- und Sensorfunktionen ausgeführt. 
Die zeitliche Steuerung erfolgt über feste Intervalle, wodurch eine gleichmäßige und deterministische Aktualisierung gewährleistet wird. 
Der vollständige Ablauf ist in Abbildung~\ref{fig:hauptprogramm} dargestellt.

\begin{figure}[H]
\centering
\begin{tikzpicture}[
    node distance=0.4cm,
    startstop/.style={rectangle, rounded corners, minimum width=4cm, minimum height=0.8cm, text centered, draw=black},
    process/.style={rectangle, minimum width=4.6cm, minimum height=0.8cm, text centered, draw=black},
    decision/.style={diamond, aspect=2.4, minimum width=4cm, minimum height=1cm, text centered, draw=black},
    arrow/.style={thick, ->, >=stealth}
]

\node (start) [startstop] {Start};
\node (setup) [process, below=of start] {setup ausführen};
\node (loop) [process, below=of setup] {loop dauerhaft wiederholen};
\node (web) [process, below=of loop] {Webserver-Anfragen bearbeiten};
\node (check500) [decision, below=of web] {500 ms vergangen?};
\node (time) [process, below=of check500] {Zeit aktualisieren};
\node (valid) [decision, below=of time] {Zeit gültig?};
\node (daily) [process, below=of valid] {Tägliche Zeitsynchronisation prüfen};
\node (weather) [process, below=of daily] {Wetteraktualisierung prüfen};
\node (sensor) [process, below=of weather] {Innenraumsensor prüfen};
\node (wordclock) [process, below=of sensor] {Wortuhr anzeigen};
\node (second) [process, below=of wordclock] {Sekunden-LED anzeigen};
\node (third) [process, below=of second] {Zusatz-LEDs anzeigen};

\draw [arrow] (start) -- (setup);
\draw [arrow] (setup) -- (loop);
\draw [arrow] (loop) -- (web);
\draw [arrow] (web) -- (check500);

\draw [arrow] (check500) -- node[right] {Ja} (time);
\draw [arrow] (time) -- (valid);
\draw [arrow] (valid) -- node[right] {Ja} (daily);
\draw [arrow] (daily) -- (weather);
\draw [arrow] (weather) -- (sensor);
\draw [arrow] (sensor) -- (wordclock);
\draw [arrow] (wordclock) -- (second);
\draw [arrow] (second) -- (third);

\draw [arrow] (check500.east) -- ++(2.0,0)
    node[midway, above] {Nein}
    |- (loop.east);

\draw [arrow] (valid.east) -- ++(2.8,0)
    node[midway, above] {Nein}
    |- (loop.east);

\draw [arrow] (third.east) -- ++(3.6,0)
    |- (loop.east);

\end{tikzpicture}
\caption{Hauptprogrammablauf der Wortuhr}
\label{fig:hauptprogramm}
\end{figure}

\subsection{Zeitverarbeitung und Synchronisation}

Die Zeitverarbeitung basiert auf der Nutzung eines \gls{ntp}-Servers zur Initialisierung sowie einer periodischen Nachsynchronisation zur Sicherstellung der Genauigkeit. 
Die aktuelle Zeit wird regelmäßig aus dem System gelesen und in globale Variablen überführt. 
Zusätzlich erfolgt täglich eine erneute Synchronisation, um Drift zu vermeiden. 
Der Ablauf ist in Abbildung~\ref{fig:zeit} dargestellt.

\begin{figure}[H]
\centering
\begin{tikzpicture}[
    node distance=1.2cm,
    startstop/.style={rectangle, rounded corners, minimum width=4cm, draw=black},
    process/.style={rectangle, minimum width=4.5cm, draw=black},
    decision/.style={diamond, aspect=2.3, draw=black},
    arrow/.style={->, thick}
]

\node (start) [startstop] {Zeit aktualisieren};
\node (read) [process, below=of start] {Systemzeit lesen};
\node (valid) [decision, below=of read] {Zeit gültig?};
\node (daily) [process, below=of valid] {Tägliche Synchronisation prüfen};
\node (end) [startstop, below=of daily] {Ende};

\draw [arrow] (start) -- (read);
\draw [arrow] (read) -- (valid);
\draw [arrow] (valid) -- node[right]{Ja} (daily);
\draw [arrow] (daily) -- (end);
\draw [arrow] (valid.west) -- ++(-2,0) node[above]{Nein} |- (end.west);

\end{tikzpicture}
\caption{Zeitverarbeitung und Synchronisation}
\label{fig:zeit}
\end{figure}

\subsection{Anzeige-Logik der Wortuhr}

Die zentrale Funktion der Software besteht in der Umwandlung der numerischen Zeit in eine sprachbasierte Darstellung. 
Dabei werden Minutenintervalle interpretiert und die entsprechende Stunde abhängig vom Kontext angepasst. 
Die Ausgabe erfolgt durch gezielte Aktivierung von \gls{led}-Gruppen, welche Wörter darstellen. 
Der Ablauf der Anzeigeerzeugung ist in Abbildung~\ref{fig:wortuhr} dargestellt.

\begin{figure}[H]
\centering
\begin{tikzpicture}[
    node distance=1.1cm,
    startstop/.style={rectangle, rounded corners, minimum width=4cm, draw=black},
    process/.style={rectangle, minimum width=4.8cm, draw=black},
    decision/.style={diamond, aspect=2.3, draw=black},
    arrow/.style={->, thick}
]

\node (start) [startstop] {Anzeige starten};
\node (clear) [process, below=of start] {LEDs zurücksetzen};
\node (power) [decision, below=of clear] {Lampe aktiv?};
\node (always) [process, below=of power] {Feste Wörter anzeigen};
\node (calc) [process, below=of always] {Anzeigestunde berechnen};
\node (minute) [process, below=of calc] {Minutenwörter anzeigen};
\node (hour) [process, below=of minute] {Stundenwort anzeigen};
\node (weather) [process, below=of hour] {Wetterinformationen anzeigen};
\node (end) [startstop, below=of weather] {Anzeige fertig};

\draw [arrow] (start) -- (clear);
\draw [arrow] (clear) -- (power);
\draw [arrow] (power) -- node[right]{Ja} (always);
\draw [arrow] (always) -- (calc);
\draw [arrow] (calc) -- (minute);
\draw [arrow] (minute) -- (hour);
\draw [arrow] (hour) -- (weather);
\draw [arrow] (weather) -- (end);

\draw [arrow] (power.west) -- ++(-2,0) node[above]{Nein} |- (end.west);

\end{tikzpicture}
\caption{Anzeige-Logik der Wortuhr}
\label{fig:wortuhr}
\end{figure}

\subsection{Erweiterung durch Sensorik und Wetterdaten}

Neben der reinen Zeitanzeige wird das System um zusätzliche Funktionen erweitert. 
Hierzu zählen die Anzeige von Wetterdaten sowie die Integration eines Innenraumsensors zur Erfassung von Temperatur und Luftfeuchtigkeit. 
Die Daten werden in definierten Zeitintervallen aktualisiert und anschließend in abstrahierter Form dargestellt. 
Der Ablauf ist in Abbildung~\ref{fig:erweiterung} dargestellt.

\begin{figure}[H]
\centering
\begin{tikzpicture}[
    node distance=1.2cm,
    startstop/.style={rectangle, rounded corners, minimum width=4cm, draw=black},
    process/.style={rectangle, minimum width=4.8cm, draw=black},
    decision/.style={diamond, aspect=2.3, draw=black},
    arrow/.style={->, thick}
]

\node (start) [startstop] {Erweiterungsfunktionen};
\node (weather) [process, below=of start] {Wetterdaten prüfen};
\node (sensor) [process, below=of weather] {Sensorwerte prüfen};
\node (display) [process, below=of sensor] {Zusatzanzeigen erzeugen};
\node (end) [startstop, below=of display] {Ende};

\draw [arrow] (start) -- (weather);
\draw [arrow] (weather) -- (sensor);
\draw [arrow] (sensor) -- (display);
\draw [arrow] (display) -- (end);

\end{tikzpicture}
\caption{Ablauf der Erweiterungsfunktionen}
\label{fig:erweiterung}
\end{figure}

\subsection{Benutzerinteraktion über Webschnittstelle}

Zur Steuerung der Wortuhr wird ein integrierter Webserver eingesetzt. 
Dieser ermöglicht die Anzeige von Statusinformationen sowie die Anpassung von Parametern wie Helligkeit, Farbe und Betriebszustand. 
Die Verarbeitung der Anfragen erfolgt ereignisbasiert und ist unabhängig vom zyklischen Hauptprogramm. 
Statusinformationen werden bei Bedarf dynamisch im \gls{json}-Format erzeugt und an den Client übertragen.
Der Ablauf ist in Abbildung~\ref{fig:webserver} dargestellt.

\begin{figure}[H]
\centering
\begin{tikzpicture}[
    node distance=1.2cm,
    startstop/.style={rectangle, rounded corners, minimum width=4cm, draw=black},
    process/.style={rectangle, minimum width=4.8cm, draw=black},
    decision/.style={diamond, aspect=2.3, draw=black},
    arrow/.style={->, thick}
]

\node (start) [startstop] {HTTP-Anfrage};
\node (path) [decision, below=of start] {Anfragepfad?};

\node (html) [process, below left=1.5cm and 1.8cm of path] {HTML-Seite erzeugen};
\node (json) [process, below=of path] {JSON-Status erzeugen};
\node (set) [process, below right=1.5cm and 1.8cm of path] {Parameter setzen};

\node (end) [startstop, below=3cm of path] {Antwort senden};

\draw [arrow] (start) -- (path);
\draw [arrow] (path) -- (html);
\draw [arrow] (path) -- (json);
\draw [arrow] (path) -- (set);

\draw [arrow] (html) |- (end);
\draw [arrow] (json) -- (end);
\draw [arrow] (set) |- (end);

\end{tikzpicture}
\caption{Ablauf der Webserver-Verarbeitung}
\label{fig:webserver}
\end{figure}